\section{Mismatch Cost and Three Regimes}\label{sec:zones}

Consider one frequency (or a band treated as a scalar channel) with signal
variance $p$ and effective noise $q$---sensor noise plus aliased energy---so
that $R=p/(p+q)$. An estimator that assumes variance $\kappa p$ instead of $p$
applies the gain $g=\kappa R/(\kappa R+1-R)$.

\begin{proposition}[Mismatch cost]\label{prop:mismatch}
The excess mean-squared error, normalised by the band energy $p$, is
\begin{equation}\label{eq:delta}
\Delta(R,\kappa)=\frac{R(1-R)^2(\kappa-1)^2}{(1-R+\kappa R)^2}.
\end{equation}
\end{proposition}
\noindent A proof is in Appendix~\ref{app:proof}. Two limits carry the
prescription:
\begin{itemize}
\item under-estimation, $\kappa\to0$: $\Delta\to R$; harmless when $R$ is low,
costly when $R$ is high;
\item over-estimation, $\kappa\to\infty$: $\Delta\to(1-R)^2/R$; harmless when
$R$ is high, and larger than the band energy itself when $R<0.38$, i.e.\ the
reconstruction is worse than not reconstructing the band at all.
\end{itemize}

\noindent In absolute units, with $Q=\kappa P$ the assumed and $P$ the true
energy and $N$ the effective noise at one frequency,
\begin{equation}\label{eq:delta-abs}
P\,\Delta=\frac{N^2(Q-P)^2}{(P+N)(Q+N)^2}.
\end{equation}
In the over-estimation window $P<N<Q$ the excess error is $\approx N$,
independent of the signal energy; in the under-estimation window $Q<N<P$ it is
$\approx P$. Over-estimation is paid in units of noise, under-estimation in
units of signal \src{DER§1.2}.

\begin{corollary}[Where the asymmetry flips]\label{cor:flip}
In the per-frequency scalar model, for every $\kappa>1$,
$\Delta(R,\kappa)>\Delta(R,1/\kappa)$ if and only if $R<1/2$.
\end{corollary}
\noindent The proof follows from the sign of $(\kappa-1)(1-2R)$
(Appendix~\ref{app:proof}). The statement is exact under two conditions: it is
evaluated per frequency, and the mismatch is a symmetric multiplicative factor
($\kappa$ vs.\ $1/\kappa$). Averaged over a band in which $R$ varies, Jensen's
inequality shifts the crossover, so band-level comparisons near $R=1/2$ are
only approximate (\S\ref{sec:exp-e2}). The model also assumes that the
effective noise $q$ does not depend on the assumed spectrum; random sparse
sampling, which aliases energy across bands, violates this, and the one
counter-example to the band-level asymmetry in E2, sparse 10\% $b_1$, lies
there (\S\ref{sec:exp-e2}).

\paragraph{Three regimes.} We fix two thresholds before the experiments,
$0.8$ and $0.1$ (sensitivity in \S\ref{sec:limits}):
\begin{itemize}
\item \textbf{Regime A} ($\Rlin\ge0.8$): the information is in the
observation. In the limit $R\to1$ over-estimating the spectrum costs
little and under-estimating it costs much, which favours wide, low-bias
representations; the bottleneck is then capacity and optimisation. This is a
per-band tendency, not a rule for choosing a representation: at $R=0.8$ a
$16\times$ over-estimate is not free, and the winner is decided by the
energy-weighted cost summed over \emph{all} bands, each with its own
$\kappa_b$ (Eq.~\eqref{eq:crit}).
\item \textbf{Regime B} ($0.1\le\Rlin<0.8$): both directions cost; the gain of
spectral adaptation is bounded by
$\mathrm{LMMSE}(\text{true})-\mathrm{LMMSE}(\text{mismatched})$, and
matching the spectrum decides the outcome.
\item \textbf{Regime C} ($\Rlin<0.1$): a linear estimator obtains almost
nothing. Under-estimation is free, over-estimation is catastrophic; only
priors that exploit cross-band statistical dependence---learned,
non-linear---can recover the band.
\end{itemize}

\paragraph{Zone labels do not pick winners.} To compare representations with
effective priors $Q$ we use the band-summed, absolute cost
\begin{equation}\label{eq:crit}
C(Q)=\sum_b E_b\,\Delta\big(R_b,\kappa_b(Q)\big),
\end{equation}
with band energy $E_b$, matched-prior recoverability $R_b$ and
$\kappa_b(Q)$ the representation's band energy relative to $E_b$; lower $C$
is predicted better. For linear arms the exact trace formula
$\mathbb{E}\|x-\hat x\|^2=\operatorname{tr}\!\big((I-H_QA)C_P(I-H_QA)^\top\big)+\sigma_n^2\|H_Q\|_F^2$
needs no approximation. Eq.~\eqref{eq:crit} ranks representations
(E2: 0.90; D11: 0.856) but does not reliably select the best among close
candidates (D11: 1/3), does not cover non-linear iterative fits, and its
a-priori version (generic spectrum for $E_b$, $\kappa_b$) has not been tested
(\S\ref{sec:exp-ab}) \src{DER§2.3}. The per-frequency form
Eq.~\eqref{eq:delta-abs} explains why over-estimation is paid in units of
noise but is not used as evidence: for non-diagonal operators (random shifts,
sparse sampling) its per-frequency predictions are inaccurate.

\paragraph{Three known failures as over-estimation.}
(i) Grid or hash fields without spectral decay assume a flat spectrum, i.e.\
$\kappa\gg1$ in the high bands, and collapse once noise lowers $R$ there;
(ii) Fourier features with large $\sigma$ in SuperF-style fitting do the same;
(iii) a regulariser of fixed strength corresponds to a fixed $\kappa$ relative
to a noise-dependent $p/q$, so its effective mismatch drifts with noise.
\S\ref{sec:exp-ab} tests these on multi-frame SR.

\begin{table*}[t]\centering\small
\caption{A-priori three-regime map: $\Rlin$ per band under the generic
$1/f^2$ prior, no data, no training \src{OUT3/d8}.
Shading: A $\ge0.8$, B, C $<0.1$.}\label{tab:map}
\begin{adjustbox}{max width=\linewidth}\begin{tabular}{llccc}\toprule
Task & $\sigma_n$ & $b_1$ & $b_2$ & $b_3$\\\midrule
Single image $\times2$ & .002 / .01 & 1.00 / 1.00 & .92 / .90 & .09 / .08\\
Single image $\times4$ & .002 / .01 & .98 / .98 & .09 / .08 & .01 / .01\\
Single image $\times8$ & .002 / .01 & .79 / .77 & .01 / .01 & .01 / .01\\
Sparse 5\% & .002 & .70 & .09 & .03\\
Sparse 10\% & .002 & .81 & .20 & .07\\
Sparse 25\% & .002 & .92 & .47 & .20\\
Sparse 50\% & .002 & .98 & .76 & .46\\
Multi-frame $\times4$, $T{=}4$ & .002$\to$.02 & $\approx$1.0 & .73$\to$.51 & .12$\to$.03\\
Multi-frame $\times4$, $T{=}8$ & .002$\to$.02 & $\approx$1.0 & .83$\to$.62 & .27$\to$.08\\
Multi-frame $\times4$, $T{=}16$ & .002$\to$.02 & $\approx$1.0 & .90$\to$.73 & .41$\to$.15\\
\bottomrule\end{tabular}\end{adjustbox}
\end{table*}

\begin{figure}[t]\centering
\fbox{\parbox{0.9\linewidth}{\centering\vspace{1.5em}Fig.~1: three-regime map
(Table~\ref{tab:map} as a heat map, four task families including EEG).
\TODO{render from d8 JSON}\vspace{1.5em}}}
\caption{A-priori regime map.}\label{fig:map}
\end{figure}

\paragraph{Two ways to assign a regime.}
$\Rlin$ describes what the observation permits a linear estimator; it is the
right axis for linear, multi-frame, sparse and 1D tasks, and it is computable
in advance. A trained ASISR network, however, lifts regime-C bands to
intermediate recovery through its learned prior (\S\ref{sec:exp-c}). To
describe what components \emph{inside} such a network do, the regime has to
be assigned by the network's actual recovery $\Rnet$, which requires ground
truth and is therefore a post-hoc description, not a prediction
(\S\ref{sec:exp-f1}).
At this network layer Prop.~\ref{prop:mismatch} is not a theorem: a network
is not a linear Wiener estimator under some prior, and has no well-defined
$\kappa$. What carries over is a boundary statement only: a component cannot
matter in bands with $\Rnet\approx1$ (the output does not depend on prior
details) or $\Rnet\approx0$ (nothing is recovered to lose), so ablation losses
can appear only at intermediate $\Rnet$. No quantitative ranking and no
direction of asymmetry follow \src{DER§3.2}.
